\begin{exercise}{40}
Suppose that $F$ is left-adjoint to $G$, witnessed by a unit and counit $(\eta,\varepsilon)$. Show
that if $F$ is monoidal, then the induced lax monoidal structure on $G$ of Exercise~\ref{ex:38} makes
$(F,G,\eta,\varepsilon)$ a monoidal adjunction.

\begin{proof}
Write $\nu,\nu_0$ for the invertible lax structure maps of
$F\colon\mathcal C\to\mathcal D$.
Their inverses $\delta=\nu^{-1}$ and $\delta_0=\nu_0^{-1}$
form an oplax structure: its naturality and coherence
diagrams follow by inverting those for $\nu,\nu_0$.
Equip $G$ with the induced lax structure $\mu,\mu_0$ of
Exercise~\ref{ex:38}.
Using \eqref{eq:monoidal-mates}, for $X,Y\in\mathcal C$ we obtain
\[
 \begin{aligned}
  G(\nu_{X,Y})\mu_{FX,FY}(\eta_X\otimes\eta_Y)
   &=G(\nu_{X,Y}\delta_{X,Y})\eta_{X\otimes Y}
    =\eta_{X\otimes Y},\\
  G(\nu_0)\mu_0
   &=G(\nu_0\delta_0)\eta_{\bbone_{\mathcal C}}
    =\eta_{\bbone_{\mathcal C}}.
 \end{aligned}
\]
For $A,B\in\mathcal D$, the other two mate identities give
\[
 \begin{aligned}
  \varepsilon_{A\otimes B}F(\mu_{A,B})\nu_{GA,GB}
   &=(\varepsilon_A\otimes\varepsilon_B)
      \delta_{GA,GB}\nu_{GA,GB}
    =\varepsilon_A\otimes\varepsilon_B,\\
  \varepsilon_{\bbone_{\mathcal D}}F(\mu_0)\nu_0
   &=\delta_0\nu_0=\mathrm{id}_{\bbone_{\mathcal D}}.
 \end{aligned}
\]
By Definition~1.2.32, the first two equations say that
$\eta\colon\mathrm{id}_{\mathcal C}\to GF$ is monoidal,
and the last two say that
$\varepsilon\colon FG\to\mathrm{id}_{\mathcal D}$ is monoidal,
using the composite lax structures on $GF$ and $FG$.
Thus $(F,G,\eta,\varepsilon)$ is a monoidal adjunction.
\end{proof}
\end{exercise}
